%%%%%%%%%%%%%%%%%%%%%%%%%%%%%%%%%%%%%%%%%%%%%%%%%%%%%%%%%%%%%%
%%%%%%%%%%%% MA-0635 Introducción a la Topología %%%%%%%%%%%%%%%%%%%
%%%%%%%%%%%%%%%%%%%%%%%%%%%%%%%%%%%%%%%%%%%%%%%%%%%%%%%%%%%%%%
\newpage
\subsection*{MA-0635 Introducción a la Topología}
\addcontentsline{toc}{subsection}{MA-0635 Introducción a la Topología}

\begin{refsection}
\begin{table}[H]
\centering{
\begin{tabular}{|ll|}
\hline
%\multicolumn{2}{|c|}{\textbf{Nivel de virtualidad del curso:} Virtual}\\ \hline
\textbf{Año:} IV  & \textbf{Requisitos:} MA-0515 \\ 
\textbf{Ciclo:} VII  & \textbf{Correquisitos:} Ninguno \\ 
\textbf{Tipo de curso:} Teórico & \textbf{Horas presenciales por semana:}   \\ 
\textbf{Créditos:} 4 & \textbf{Horas de trabajo independiente por semana:}  \\ \hline
\end{tabular}
}
\end{table}



\subsubsection*{Descripción del curso}

Este es un curso del sexto ciclo de la carrera Bach. y Lic. en Matemática en el que se presentan las bases de la topología general. El curso proporciona los cimientos para el estudio de propiedades de los espacios topológicos y las funciones entre ellos, a través de una perspectiva rigurosa. Se inicia con la exploración de espacios métricos, lo que permite extender de manera intuitiva a estos espacios más generales, conceptos centrales como continuidad, conjuntos abiertos y cerrados, compacidad y conexidad, que han sido previamente estudiados en el contexto de $\R^n$. A lo largo del curso, se profundiza en estos conceptos, abordando también los axiomas de separación, topologías producto, subespacios, el grupo fundamental y homotopía, que son herramientas esenciales en el análisis y comprensión de la estructura de diferentes tipos de espacios.

La topología es una disciplina fundamental en la matemática moderna, ya que ofrece una estructura unificadora para el análisis en diversos contextos. Su relevancia se extiende desde el análisis real y complejo hasta áreas como la geometría diferencial, la teoría de la medida y la teoría de números. Además, los conceptos topológicos son ampliamente aplicables en ciencias aplicadas y ramas emergentes como la topología de datos y la teoría de redes. Con una sólida base en topología general, el estudiante estará preparado para abordar problemas en matemáticas avanzadas y comprender la estructura profunda de los objetos matemáticos.



\subsubsection*{Objetivos}

\textbf{General:} Aplicar los conceptos y resultados básicos de la topología general para la resolución de problemas teóricos y prácticos. \\

\textbf{Específicos:} Al finalizar el curso se espera que la persona estudiante sea capaz de:

\begin{enumerate}
\item Comparar los conceptos, propiedades e implicaciones de la convergencia, continuidad, diferenciación e integración de funciones de variable real y funciones de variable compleja.
\item Demostrar resultados sobre límites, continuidad, convergencia, diferenciación e integración en el análisis de funciones de una variable compleja.
\item Aplicar los teoremas de diferenciación e integración de funciones de variable compleja en la demostración de resultados y la resolución de problemas teóricos y prácticos.
\item Relacionar aspectos algebraicos, analíticos y geométricos de los números complejos y de las funciones de variable compleja.
\item Indagar en la literatura sobre temas que permitan complementar su formación.
\item Comunicar ideas matemáticas de manera clara y efectiva en forma oral y escrita.
\end{enumerate}

\subsubsection*{Contenidos}



\begin{enumerate}

\item \textbf{Repaso de teoría de conjuntos (1 semana):} 
\begin{enumerate}
    \item Definiciones y operaciones básicas.
    \item Familias de conjuntos: uniones e intersecciones.
    \item Producto cartesiano de una familia de conjuntos: axioma de elección.
    \item Funciones y cardinalidad: ejemplos y resultados básicos.
\end{enumerate}

\item \textbf{Espacios métricos (3 semanas):} 
\begin{enumerate}
    \item Espacios métricos: definiciones, ejemplos y construcciones básicas (subespacios y productos de espacios métricos).
    \item Bolas, conjuntos abiertos y cerrados, clausura de un subconjunto, subconjuntos densos.
    \item Sucesiones y convergencia. Completitud y compleción de un espacio métrico. 
    \item Compacidad: definiciones, ejemplos, el Teorema de Heine-Borel.
    \item Continuidad: definiciones, ejemplos, continuidad uniforme.
    \item Isometrías y homeomorfismos: definiciones, ejemplos y equivalencias de métricas.
\end{enumerate}

\item \textbf{Espacios topológicos (4 semanas):} 
\begin{enumerate}
    \item 
\end{enumerate}

\item \textbf{Axiomas de separación (3 semanas):} 
\begin{enumerate}
    \item 
\end{enumerate}

\item \textbf{Homotopía y el grupo fundamental (3 semanas):} 
\begin{enumerate}
    \item 
\end{enumerate}

\item \textbf{Tema optativo (1 semana):} 
\begin{enumerate}
    \item 
\end{enumerate}


\end{enumerate}

\subsubsection*{Metodología}

\noindent \textbf{Lineamientos metodológicos:} La persona docente a cargo de este curso debe respetar las siguientes pautas:
\begin{enumerate}
    \item Fomentar el desarrollo de intuición sobre los conceptos estudiados en el curso.
    \item Promover la comunicación clara y rigurosa de argumentos e ideas matemáticas de manera oral y escrita.
    \item Cuando la naturaleza del contenido lo permita, se deben utilizar representaciones visuales que apoyen la comprensión de los temas.
    \item Fomentar una visión integral de las matemáticas y de cómo se enlazan las diferentes áreas a través de la topología.
\end{enumerate}

\textbf{Sugerencias metodológicas:} Adicionalmente, se tienen las siguientes recomendaciones para la persona docente a cargo de este curso:
\begin{enumerate}
    \item Utilizar el recurso tecnológico para reforzar distintos conceptos estudiados en el curso por medio de la visualización.
    \item Aprovechar momentos adecuados del curso para motivar el estudio por medio de reseñas acerca del desarrollo histórico de la topología y su importancia en aplicaciones dentro y fuera de la matemática.
    \item En la evaluación, se sugiere utilizar estrategias que promuevan el trabajo constante de parte de las personas estudiantes a lo largo del ciclo lectivo. Por ejemplo, esto se puede hacer por medio de listas de ejercicios recomendados o proyectos.
    \item Para fomentar el desarrollo de las habilidades investigativas en el estudiantado, se sugiere la implementación de proyectos de investigación y participación en coloquios o seminarios de investigación.
    \item Dado que hoy en día la mayoría del trabajo en matemática, tanto a nivel académico como profesional, es de naturaleza colaborativa, se sugiere a la persona docente implementar estrategias que promuevan el trabajo en equipo.
\end{enumerate}



\nocite{Apo76}
\nocite{Mun00}
\nocite{GG99}
\nocite{}
\nocite{}




\printbibliography[heading=subbibliography]
\end{refsection}
